\documentclass[11pt,letterpaper]{article}
\usepackage[T1]{fontenc}
\usepackage[utf8]{inputenc}
\usepackage{lmodern,microtype,amsmath,amssymb,amsthm,mathtools}
\usepackage[margin=1in]{geometry}
\usepackage{booktabs,longtable,tabularx,array,enumitem,xcolor,listings,fancyhdr}
\usepackage[colorlinks=true,linkcolor=blue!45!black,urlcolor=blue!50!black,citecolor=blue!45!black]{hyperref}
\usepackage{xurl}
\hypersetup{pdftitle={K215: Direct Seed Models and Shared Propositional Proofs},pdfauthor={OpenAI, research assistance for Vladimir Reshetnikov}}
\setlist{itemsep=3pt,topsep=6pt}
\setlength{\emergencystretch}{2em}
\setlength{\headheight}{14pt}
\pagestyle{fancy}
\fancyhf{}
\lhead{\small ProveIt / K215 proof simplification}
\rhead{\small October 21, 2026}
\cfoot{\thepage}
\newtheorem{theorem}{Theorem}[section]
\newtheorem{lemma}[theorem]{Lemma}
\newtheorem{proposition}[theorem]{Proposition}
\theoremstyle{definition}
\newtheorem{definition}[theorem]{Definition}
\theoremstyle{remark}
\newtheorem{remark}[theorem]{Remark}
\newcommand{\enc}{\operatorname{enc}}
\newcommand{\thh}{\operatorname{th}}
\newcommand{\Sat}{\operatorname{Sat}}
\newcommand{\Free}{\operatorname{Free}}
\newcommand{\Seed}{\operatorname{Seed}}
\newcommand{\Der}{\operatorname{DerivableNF}}
\newcommand{\FullSeed}{\operatorname{FullSeed}}
\newcommand{\Ccal}{\mathsf{C}_{222}}
\newcommand{\NF}{\mathsf{NF}}
\newcommand{\code}[1]{\texttt{\detokenize{#1}}}
\newcommand{\pin}{\texttt{6d9800d1ca67c52db0eb6e1c65df7b5cab9cdb95}}
\newcommand{\project}{\url{https://github.com/VladimirReshetnikov/ProveIt}}
\lstset{basicstyle=\ttfamily\small,columns=fullflexible,keepspaces=true,breaklines=true,
  frame=single,rulecolor=\color{black!15},backgroundcolor=\color{black!2},
  xleftmargin=4pt,xrightmargin=4pt,aboveskip=10pt,belowskip=10pt,showstringspaces=false}
% Values are bound to the final evidence files in this packet.
\newcommand{\LeanBeforeModules}{178}
\newcommand{\LeanAfterModules}{177}
\newcommand{\LeanBeforeBytes}{2,093,112}
\newcommand{\LeanAfterBytes}{2,086,255}
\newcommand{\LeanBytesSaved}{6,857}
\newcommand{\LeanLinesSaved}{130}
\newcommand{\RocqBytesSaved}{7,159}
\newcommand{\RocqLinesSaved}{36}


\title{\textbf{K215: Direct Seed Models and\\Shared Propositional Proofs}\\[5pt]
\large Two checked components for simplifying the Rocq and Lean developments}
\author{Research assistance for Vladimir Reshetnikov\\\small OpenAI}
\date{October 21, 2026 (America/Los\_Angeles)}

\begin{document}
\maketitle
\begin{abstract}
The current K215 development already contains many of the obvious proof reductions.
This report develops two specific next steps against the pinned repository: a direct
Lean construction of a full seed model, and a shared Rocq propositional proof core.
The direct construction implements an earlier uncompiled repository proposal and
aligns Lean with a construction already used by the Rocq lower bound. Its essential
point is that the completeness interface supplies truth of original named formulas
with literal equality. One can therefore build the seed model on the original
domain, using an interpreted formula code as a fresh private name, and bypass the
preliminary indiscernibility quotient. The Rocq refactor packages 71 duplicated
derivations once over implication, negation, modus ponens, and the three existing
propositional axiom rules. Concrete adapters preserve both sets of exported names.
The archive contains source overlays, reproduction programs, exact input identities,
compiler evidence, and independent review. The status and limits of compilation are
stated separately for the semantic components and the complete physical-machine
endpoint. No Turing-machine transition is changed.
\end{abstract}

\tableofcontents
\newpage

\section{Recommendation and scope}

The recommended next integration is \emph{direct seed construction in Lean plus
one shared propositional core in Rocq}. These remove different kinds of unnecessary
work: an intermediate model construction in the first case, and duplicate derivation
bodies in the second. They preserve the target theorem and the existing calculus.
They do not search for a different machine or rely on a smaller machine to simplify
the proof.

\subsection{Precisely which repository state was reviewed}

The repository is \project. All source comparisons in this report use commit
\pin, whose latest relevant change generalizes the K217 enumerator to arbitrary
states between candidates. The timestamp is October 22 in UTC and October 21 in
the user's local time zone. The previous changes already include the phase-fourteen
K215 checkpoint and the proposed CP/GIC-to-QHyp quotient adapter.

The independent source census reproduces the phase-fourteen Lean endpoint exactly:
\textbf{178 project modules, 2,093,112 source bytes, and 40,517 physical source
lines}. This is an import-closure measurement; it is not a count of mathematical
lemmas or a measurement of proof-object size. The source identity and complete
module list are retained in \code{evidence/baseline-lean-closure.json}.

The dense machine table has SHA-256
\begin{quote}\small\ttfamily
fdbcef0825886035c3091c8537a69535bc96f9e265f0ae45863543d40ae70090.
\end{quote}
It has 215 working states and a separate halt. Neither proposal changes that table.

\subsection{The theorem being preserved}

The maintained Lean endpoint has type
\begin{lstlisting}
BlankHalts machine215 <-> BProv ZFCax_s [] Form.fBot
\end{lstlisting}
and its nonhalting corollary negates both sides. The Rocq endpoint similarly
identifies halting of the exact dense table with a finite classical refutation of
its explicitly defined ZFC theory. These are external proof-assistant theorems
about the machine and syntactic consistency. The repository also explains why the
ordinary mathematical argument can be formalized in ZFC. The present changes do
not supply a serialized derivation in ZFC's own proof syntax, and they do not close
the separately open internal PA endpoint.

The distinction matters: proving a new component, checking a source overlay, and
rebuilding the entire physical endpoint are three different achievements. We use
each description only for its corresponding evidence.

\subsection{Measured result and verification boundary}

\begin{center}
\begin{tabularx}{\textwidth}{@{}p{0.28\textwidth}p{0.20\textwidth}p{0.20\textwidth}X@{}}
\toprule
Measured source scope & Before & After & Change \\
\midrule
Lean K215 import closure & 178 modules & 177 modules & One fewer module \\
Same closure, bytes & 2,093,112 & 2,086,255 & 6,857 fewer \\
Same closure, physical lines & 40,517 & 40,387 & 130 fewer \\
Rocq paired propositional sources, bytes & 23,160 & 16,001 & 7,159 fewer \\
Same sources, physical lines & 297 & 261 & 36 fewer \\
Maintained Rocq derivation bodies & 142 & 71 & One copy per lemma \\
\bottomrule
\end{tabularx}
\end{center}

The Lean figure measures the \emph{active endpoint closure}. Keeping the old
quotient interfaces available adds 7,472 bytes and 184 lines across the six
changed or new production files overall; compatibility preservation therefore
does not shrink the entire repository. The three modules leaving this endpoint
closure are \code{Code254Models}, \code{KSeedModel222}, and \code{KIndisc222};
two small interface/direct-proof modules enter it. The numerical seed semantics
module remains because the upper chain still uses it.

The checked scope is the new Lean semantic component and its recorded clients,
and the Rocq propositional component with its actual consumers. The entire
literal K215 endpoint has not been rebuilt in this packet. The repository's
existing full-endpoint checks remain evidence for the original pinned sources;
the changed endpoint requires the stated integration gate.


The reductions are intentionally reported with their costs. Adding a shared
abstraction is useful only when its consumers actually share it. Likewise,
removing a quotient from an informal diagram does not remove it from the
compiled dependency closure until imports and shared declarations are separated.

\section{What recent work has already accomplished}

This work begins after substantial simplification. Treating earlier results as new
would obscure the remaining problems. The current source and proposal index establish
the following context; references are given in Section~\ref{sec:references}.

\begin{longtable}{@{}p{0.28\textwidth}p{0.67\textwidth}@{}}
\toprule
Existing work & Consequence for the present recommendation \\
\midrule
\endhead
Shared implication core and named derived rules & The numerical dominated
implication argument and the two Lean named calculi already share major parts of
their derivations. The Rocq target below is a separate pair of classical
propositional proof files. \\[4pt]
Transitive-container factoring & Phase fourteen removes a duplicate profile
development and its generator. We use that checkpoint as the baseline. \\[4pt]
Checked macro words & The Lean AX construction for the permuted source is already
a checked word; the dominated source now uses the same template. Recommending
that transformation again would add little. \\[4pt]
Source and runtime certificates & The proposals already cover trailing-zero
normalization, balanced lookup, affine cores, semantic site contracts, uniform
decoder prefixes, and positive simulations. They retain their own adoption and
verification status. \\[4pt]
Candidate scheduling & A generic scheduling induction is already proposed, and
the current K217 change removes a seed-specific repetition. \\[4pt]
Quotient unification & A recent proposal derives the QHyp interface from CP/GIC
and reuses one later bisimulation-quotient development. That proposal is distinct
from removing the \emph{earlier} indiscernibility quotient. \\[4pt]
Direct full seed & The sharing-interfaces packet explicitly sketches the direct
Lean route but labels it uncompiled. The corresponding Rocq construction already
appears in \code{rocq224/src/Complete222.v}. The contribution here is checked
implementation and integration work, with those origins acknowledged. \\
\bottomrule
\end{longtable}

The common design principle is to expose the strongest useful theorem that is
already available, and then prove each reusable argument at the weakest interface
it actually needs. The direct-seed result uses a stronger semantic interface than
the old caller exploited. The propositional core uses a weaker syntactic interface
than its old proof files suggested.

\section{The Lean opportunity: use the original model}

\subsection{The current lower chain}

The numerical calculus \(\Ccal\) derives the code \(1\) exactly when the relevant
set theory is refutable. The difficult lower direction starts with a hypothetical
model of the named formulas whose equality-interpreted numerical codes are
derivable. In the previous route, the named seed bundle was first obtained only
for equality-free matrices. The model was then quotiented by membership
indiscernibility so that its new equality could be represented by an equality-free
formula upstairs. A satisfaction-transport induction and representative choices
extended the seed property to all named matrices.

This is a valid general construction. However, the actual completeness interface
already supplies more than an arbitrary model of equality-free formulas.

\subsection{The decisive completeness interface}

Write \(\thh f\) for the membership-indiscernibility interpretation of equality,
and \(\enc f\) for the numerical encoding of a named formula. The theory used in
the lower bridge is
\[
  \Der(f) \quad\Longleftrightarrow\quad
  \Ccal\bigl(\enc(\thh f)\bigr).
\]
The syntactic premise of \code{KEind222.c222_complete} asks for derivability of
\(\enc(\thh f)\) for theory members \(f\). Its semantic premise, however, supplies
\[
  \forall f\;\bigl(\Der(f)\Longrightarrow
       \forall e\;\Sat(M,e,f)\bigr).                 \tag{1}
\]
It supplies truth of \(f\), with literal equality, on the original domain of \(M\).
It does not merely supply truth of \(\thh f\). This is the fact that permits the
shorter route.

There is no circularity. The completeness theorem is proved using the named
calculus and its interpretation; it does not import the seed-model or
indiscernibility developments that this refactor bypasses.

\subsection{The seed property}

Let \(E\) be an arbitrary binary relation on a domain \(V\). Write \(x\subseteq_E y\)
for \(\forall w\,(E(w,x)\to E(w,y))\). For a binary predicate \(Q\) and an input
\(A\), the seed property is
\begin{align*}
\Seed_E(Q,A)\;\Longleftrightarrow\;\exists C\;\bigl[&
  (\forall y\,(Q(A,y)\to\exists z\,(Q(A,z)\land E(z,C))))\\
  {}\land{}&\forall y\,(E(y,C)\to\exists x\,[Q(x,y)\land\\[-2pt]
  &\hspace{22mm}\forall x'\,(x'\subseteq_E y\to
       \forall y'\,(Q(x',y')\to\exists z\,(Q(x',z)\land E(z,C))))])\bigr].
\end{align*}
Thus \(C\) answers any output request of \(A\), and each element of \(C\) is an
output whose subsets also have their requests answered inside \(C\).

A full seed model is irreflexive and has this property for every named formula
\(\psi\), with the relation
\[
  Q_{\psi,e}(x,y)=\Sat(E,e[0\mapsto x][1\mapsto y],\psi).
\]
Equality is permitted in \(\psi\), and all other free names are interpreted as
parameters by \(e\).

\section{A direct full-seed theorem}

\subsection{A fresh name supplied by the interpreted code}

The repository's numerical constructors are
\[
\begin{array}{ll}
  \enc(R(i,j))=4P(i,j)+1,&\enc(i=j)=4P(i,j)+3,\\
  \enc(\bot)=1,&\enc(f\to g)=4P(\enc f,\enc g),\\
  \enc(\forall i\,f)=4P(i,\enc f)+2,
\end{array}
\]
where \(P\) is the pairing function already used by the machine development.
The existing projection bounds give \(x\le P(x,y)\) and \(y\le P(x,y)\).

\begin{lemma}[Free names are below the code]\label{lem:free-code}
For every named formula \(f\) and natural number \(i\),
\(\Free(i,f)\) implies \(i<\enc f\).
\end{lemma}
\begin{proof}
Induct on the formula. For membership and equality atoms, \(i\) is one of the
two names, and the pairing projection bounds together with the positive tag give
the inequality. Falsity has no free names. In an implication, apply the induction
hypothesis to the subformula in which \(i\) is free, and then the corresponding
pairing projection bound. In a universal formula, free occurrence implies free
occurrence in its body, and the body-code projection bound applies.
These are the five branches of \code{free_lt_enc}.
\end{proof}

\begin{lemma}[Equality interpretation preserves free support]\label{lem:free-th}
If \(\Free(i,f)\), then \(\Free(i,\thh f)\).
\end{lemma}
\begin{proof}
The relation, falsity, implication, and universal cases follow directly from the
definition and induction. Equality \(a=b\) is interpreted by the conjunction of
agreement of incoming and outgoing membership. Its auxiliary quantified name is
\(a+b+1\), which differs from both \(a\) and \(b\); consequently both original
argument names remain free. This proves the equality case. The assertion concerns
\emph{free} occurrences, not every syntactic occurrence.
\end{proof}

For an arbitrary named matrix \(\psi\), including equality, set
\[
  m_\psi=\neg\neg\psi,\qquad
  k_\psi=\enc(\thh m_\psi).
\]
The double negation gives \(k_\psi>1\) directly from the outer implication code
and the right projection bound \(1\le P(\enc(\thh\neg\psi),1)\).
Moreover, \(k_\psi\) is not free in \(m_\psi\). If it were free, the two lemmas
would give
\[
  k_\psi<\enc(\thh m_\psi)=k_\psi,
\]
a contradiction. Thus one and the same code supplies the enabled numerical
bundle label and the freshness needed for the named semantic proof. No search for
an extra fresh variable or alpha-renaming pass is required.

\subsection{Named bundle semantics}

The named syntax has genuine falsity, and its negation is implication into that
falsity. Let \(B(a,b,k;m)\) denote the named seed bundle \code{bundleNF}. Its
request and closure subformulas are the existing \code{reqNF} and
\code{closNF}. The exact names \(a,b,k\) remain explicit.

\begin{theorem}[Named bundle equivalence]\label{thm:bundle}
Suppose \(a\ne b\), \(k\ne a\), \(k\ne b\), and \(k\) is not free in \(m\).
Then, for every relation \(E\) and environment \(e\),
\[
  \Sat(E,e,B(a,b,k;m))\quad\Longleftrightarrow\quad
  \Seed_E\bigl((x,y)\mapsto\Sat(E,e[a\mapsto x][b\mapsto y],m),e(a)\bigr).
\]
No irreflexivity hypothesis is required for this named-formula theorem.
\end{theorem}
\begin{proof}
First normalize the named request. Since \(k\ne b\), it says exactly that every
output of the fixed input has some output in the container named by \(k\).
Repeated assignments to \(b\) collapse by the environment-update identity.

Next normalize the closure subformula. Commuting updates at the distinct names
\(a,b,k\) gives: every member \(y\) of the container is an output of some input
\(x\), and every subset \(x'\) of \(y\) has each output request answered in the
container. This is \code{named_clos_iff}.

The bundle's auxiliary implication from a membership atom to its
\(k\)-universal closure is valid because \(k\) differs from its two names.
The outer negated universal implication is then converted, classically, to an
existential container satisfying the closure and request conditions. Finally,
freshness of \(k\) in \(m\) removes the container assignment from every
occurrence of the matrix. The resulting conjunction is the definition of
\(\Seed_E\), in reversed order.
\end{proof}

\begin{remark}[Why the numerical theorem still needs irreflexivity]
The decoded numerical code \(1\) is represented by a self-membership atom;
irreflexivity makes it false. Named \(\NF.\bot\), on the other hand, is interpreted
as literal falsehood. Removing irreflexivity from the named theorem does not
permit removing it from the numerical soundness theorem. In particular, the
existing soundness development still covers arbitrary natural codes and
coincident input/output roles.
\end{remark}

\subsection{Derivability of the named bundle}

Equality interpretation commutes with the named bundle:
\[
  \thh B(a,b,k;m)=B(a,b,k;\thh m).
\]
Every atom introduced by the bundle macros is a membership atom, so this identity
requires no freshness side condition. With \(k=k_\psi\), the existing encoding
identity therefore gives
\[
 \enc\bigl(\thh B(0,1,k_\psi;m_\psi)\bigr)
   =\code{bundle222}(0,1,k_\psi).
\]
Since \(k_\psi>1\), the numerical bundle constructor proves the right side
derivable. Hence \(B(0,1,k_\psi;m_\psi)\) belongs to \(\Der\).

\begin{theorem}[Direct full seed]\label{thm:direct}
Let \(E\) be an irreflexive relation satisfying every \(\Der\) formula under
every environment. Then \(\FullSeed(E)\) holds on its original domain.
\end{theorem}
\begin{proof}
For arbitrary \(\psi,e\), apply the model hypothesis to the derivable named
bundle just constructed. Theorem~\ref{thm:bundle}, the bound \(k_\psi>1\), and
freshness give \(\Seed_E(Q_{m_\psi,e},e(0))\). Classical double-negation
elimination gives pointwise equivalence of \(Q_{m_\psi,e}\) and
\(Q_{\psi,e}\). The existing \code{seedProp_congr} transports the seed property
across this pointwise equivalence. Combining these instances with the supplied
irreflexivity gives \(\FullSeed(E)\).
\end{proof}

This is precisely the new direct semantic component. The older theorem concludes
\(\FullSeed(\operatorname{qmem}E)\) on a quotient domain; the direct theorem
concludes \(\FullSeed(E)\).
The argument does not assert that \(\psi\) and \(\thh\psi\) have the same
truth value on the original domain. The interpreted formula is used to prove
derivability and choose the private name; the model hypothesis evaluates the
original named bundle.

\section{Integration into the Lean lower bridge}

\subsection{What can now be removed from this proof path}

The old path needs the relation \code{Ind}, its setoid, the quotient
\code{IndQ}, lifted membership \code{qmem}, equality and surjectivity facts for
representatives, a satisfaction-transport induction, and a transfer of seed
containers through quotient representatives. None of these is needed by the
direct lower bridge.

The later set-theoretic construction still needs its tagged domain, bisimulation
quotient, and well-founded part. Those solve a different mathematical problem.
The proposed CP/GIC-to-QHyp adapter can be considered separately; this report's
direct seed result neither assumes its adoption nor claims its savings.

\begin{center}
\begin{tabularx}{\textwidth}{@{}p{0.18\textwidth}XX@{}}
\toprule
Stage & Previous lower bridge & Direct lower bridge \\
\midrule
Completeness & Model of original derivable named formulas & Same model \\
Seed instances & Equality-free matrices first & All named matrices directly \\
First quotient & Membership indiscernibility and satisfaction transport & Not used \\
Seed model & Quotient domain & Original domain \\
ZF construction & Tagged domain, bisimulation quotient, well-founded part & Same construction \\
Witness & Quotient class of \(e(0)\) & \(e(0)\) \\
\bottomrule
\end{tabularx}
\end{center}

\subsection{A small interface is essential}

Simply importing the old quotient module to obtain its \code{FullSeed} declaration
would retain the entire old route. The implementation therefore separates the
shared named bundle definitions, the theory predicate, and the full-seed
specification from the optional quotient proof. The downstream
\code{KLowerSeed} module needs only this specification. Existing quotient clients
can continue to use the legacy module and its original theorem statements.

This is a compatibility-preserving import cut, not deletion of a useful general
quotient theorem from the repository. The exact overlay and source census record
which modules leave the K215 closure and which declarations remain available to
other clients.

\subsection{The endpoint substitution}

At the mathematical boundary, the lower bridge changes to the following shape:
\begin{lstlisting}
have hfull := fullSeed_direct_of_derivable mem hirr hT
exact zf_consistent_of_formSeed
  (formSeed_of_fullSeed hfull) hfull.irr (e 0) hf
\end{lstlisting}
The source calculus, its code \(1\), the definition of finite ZF refutation,
the ZF/ZFC refutation equivalence, and the physical-machine theorem type are
unchanged. In particular, this argument does not add an inhabited-domain axiom:
the environment already supplies the witness \(e(0)\).

Removing a semantic quotient does not imply removing \code{Quot.sound} from
Lean's printed axiom set. Existing classical reasoning and extensionality
infrastructure still use Lean's standard principles. The relevant check is that
the exact reported axiom profile is preserved and no \code{sorryAx}, new axiom,
or unchecked native-computation principle enters.

\section{The Rocq opportunity: one propositional derivation library}

\subsection{The duplicated fragment}

The pinned \code{ZF.L1_LProp} file and the K215-local \code{L224Prop} file
contain 71 corresponding lemmas: 59 intermediate derivations and 12 public
propositional rules. An independent normalization check finds identical statements
and proof bodies after renaming the theorem predicate, the axiom constructors,
modus ponens, and the exported prefixes.

The two full named calculi should not be identified. They have different
quantifier/vacuity interfaces. The repeated \emph{propositional} derivations,
however, use only implication, negation, modus ponens, and three rules common
to both calculi.

\subsection{The minimal interface}

Let \(F\) be any type of formula objects, with operations \(\Rightarrow\) and
\(\neg\), and predicate \(\operatorname{Pr}\). Require
\begin{align*}
\mathrm{MP}:\quad&\operatorname{Pr}(a\Rightarrow b)\to
                    \operatorname{Pr}(a)\to\operatorname{Pr}(b),\\
\mathrm{B1}:\quad&\operatorname{Pr}\bigl((a\Rightarrow b)\Rightarrow
                   ((b\Rightarrow c)\Rightarrow(a\Rightarrow c))\bigr),\\
\mathrm{B2}:\quad&\operatorname{Pr}\bigl((\neg a\Rightarrow a)\Rightarrow a\bigr),\\
\mathrm{B3}:\quad&\operatorname{Pr}\bigl(a\Rightarrow(\neg a\Rightarrow b)\bigr).
\end{align*}
The Rocq module type \code{PropositionalBasis} expresses exactly these data.
The functor \code{PropositionalRules} contains one copy of the inherited
derivations. No syntax induction, quantifier constructor, equality atom, or
set-theoretic axiom is part of this contract.

Three original proof bodies used the arbitrary formula \code{Weq 0 0}, although
their statements did not mention it. Each lemma already binds a formula
\code{w0}; substituting that formula for the arbitrary choice removes even the
need for a distinguished formula or an inhabited-carrier assumption. The core
therefore applies uniformly to an arbitrary formula type, including an empty
one.

\subsection{Why the abstraction is sound}

Every original derivation is a finite term constructed from earlier derivations,
MP, and the three basis rules. Replacing the concrete symbols with interface
fields yields the same well-typed proof term over arbitrary interface data.
This is the mathematical content of the abstraction. Kernel compilation checks
each such term; the code does not assume a broad soundness principle for the
full named calculus.

For the ordinary calculus, instantiate the interface using
\code{lprov}, \code{LDET}, and \code{LB1}--\code{LB3}. For the restricted
calculus, instantiate it using \code{lprov224}, \code{MDET}, and
\code{MB1}--\code{MB3}. Both use their original implication and negation
constructors. The compatibility facades retain the original public and helper
names.

The module-type fields are parameters of a functor, not new global axioms of
the concrete endpoints. Their implementations are the constructors already
present in the original inductive derivability relations. The concrete
assumption reports must therefore be closed; this is checked explicitly.

\subsection{The real size tradeoff}

For the two concrete clients together, the old source occupies
\textbf{297 lines and 23,160 bytes}. The shared core plus its two compatibility
facades occupies \textbf{261 lines and 16,001 bytes}: a reduction of
\textbf{36 lines and 7,159 bytes}, approximately \textbf{30.9\%} of those source
bytes. These figures include the new abstraction and both facades, but exclude
audit clients and documentation from both sides of the implementation comparison.

Adopting only the restricted-calculus facade while retaining the complete old
ordinary propositional development would \emph{increase} the selected source by
97 lines and 3,653 bytes. The paired adoption is the actual simplification.
The report does not infer a proof-object-size reduction or a clean-build speedup
from these source figures; those are separate measurements.

\section{Verification, reproduction, and adoption}

\subsection{Evidence is tied to source identities}

The archive includes the repository pin, source SHA-256 values, the proposed
overlays, compiler logs, assumption reports, typed API checks, and source metrics.
The Lean closure census is independent of the compiler and labels its output as
a source measurement. The Rocq evidence records the selected component and its
upstream references. No old proof receipt is treated as a certificate for newly
combined source bytes.

\subsection{Lean verification}

The proof compiler is Lean 4.32.0, commit
\code{8c9756b28d64dab099da31a4c09229a9e6a2ef35}, matching the repository's
toolchain. The direct bundle theorem, the free-name bound, preservation of free
support through equality interpretation, and the direct seed theorem have been
compiled against the reviewed source. The successful named-bundle and direct
seed reports use only \code{propext}, \code{Classical.choice}, and
\code{Quot.sound}.

The final clean replay compiles 112 project modules serially with no project
object reuse. Every compilation succeeds. The two typed audit clients check
the generic endpoint and both direct/legacy seed APIs, producing eleven
expected axiom reports. The two free-support lemmas use \code{propext} and
\code{Quot.sound}; the other nine audited results use the standard three
axioms. There are no warnings in the changed modules or audit clients; 18
preexisting linter warnings occur in unchanged dependencies. The recorded
elapsed time is 113.52 seconds for this replay, with no baseline speed
comparison implied.

The exact compilation records are described in \code{lean/README.md} and
its accompanying evidence. The interface split is
checked separately from the optional old quotient route, so retaining old
clients cannot silently reintroduce that route into the new K215 import closure.
The source-only census computes the full hypothetical K215 closure of the
overlay, including its unchanged physical certificate modules; it does not
compile those certificates.

\subsection{Rocq component verification}

The local component check uses Coq 8.18.0 with OCaml 4.14.1. Ten selected source
modules are compiled serially, including the shared core, both actual
derivability instances, both compatibility facades, and the unchanged
\code{ZF.L1_LMeta} and \code{K215F.L224Meta} consumers. The audit restates all
142 original theorem signatures. Its 26 printed assumption reports are closed:
12 public propositional rules in each calculus and two meta-theory results.
Recursive \code{coqchk} succeeds for the selected consumers and their imports.

The only compatibility edit in the old meta-theory consumer for this local
check is the import-prefix spelling \code{Stdlib} to \code{Coq}. The packet
does not claim that Coq 8.18 is the production toolchain: the maintained K215
Rocq lane pins Rocq 9.1.0 and stdlib 9.0.0. The supplied component must be
registered in that driver and checked there before adoption. Raw exit codes,
source and object hashes, API logs, and the recursive kernel-check record are
retained in \code{rocq/evidence/}.

\subsection{Independent review and limits of performance claims}

An independent source-and-receipt review verifies the direct argument's
completeness premise, the exact import counts, the moved interface statements,
the 142 Rocq API signatures, ten Rocq source/object records, and the closed
assumption reports. The final Lean replay has its own 112 bound module
records. This review is labeled as a review, not a second compiler
run. Its detailed findings are included with the archive.

The evidence contains per-module elapsed times to reproduce the work's scope.
These are not a controlled before/after build benchmark. No clean-build
speedup or proof-object-size reduction is claimed. Module counts, physical
source lines, source bytes, and number of maintained derivation bodies are
kept as separate metrics.


\subsection{Suggested integration order}

\begin{enumerate}
\item Apply the Lean interface split and direct-seed component to the pinned
      source in an isolated checkout. Compile the direct theorem, its semantic
      lower bridge, and the retained legacy interfaces. Compare exact theorem
      types and printed axioms.
\item Run the repository's K215 endpoint verification driver on the changed
      closure. Use its exact-source provider reuse policy for unchanged physical
      modules; do not relabel a source census as a successful machine build.
\item Adopt the Rocq core and both facades together in the source-staging
      pipeline. Compile the exact production namespace and consumers under its
      pinned Rocq/stdlib versions, and run the corresponding recursive kernel
      check.
\item Update generators or staging rules so regeneration reproduces the
      simplified sources. Retain the derivative-file attribution and upstream
      pins. A locally edited generated file alone is not a durable integration.
\item Record separate results for logical API preservation, source reduction,
      proof objects, and build time. Only then replace the currently approved
      endpoint provider records.
\end{enumerate}

The work delivered here is a proposal package for this sequence. It does not
modify the public repository or ask for permission to publish changes. The
user requested an archive for later incorporation.

\section{What should not be simplified away}

Several tempting reductions would change the theorem rather than simplify its
proof.

\paragraph{Arbitrary numerical codes.}
The machine manipulates arbitrary natural formula codes. Restricting a numerical
soundness theorem to neatly decoded formulas would require a new reachability
invariant. The direct lower bridge has no such consequence for the upper bridge.

\paragraph{Coincident role names.}
The ordinary named-bundle theorem uses distinct input and output names. The
numerical soundness theorem separately handles degenerate roles. Its diagonal
case must remain; the direct construction simply selects the fixed distinct
roles 0 and 1 for its own use.

\paragraph{The later set-theoretic quotient.}
The preliminary indiscernibility quotient enforces an equality interpretation.
The later bisimulation quotient belongs to the ZF model construction. Removing
the first does not justify removing the second.

\paragraph{Progress in machine simulations.}
The present proposal does not alter the physical runtime simulation. Existing
positive-step obligations prevent a target halt from being hidden by an infinite
sequence of zero-duration source stutters. A future runtime compression must
preserve those obligations or supply an independently proved well-founded rank.

\paragraph{Exact statements and assumptions.}
A shorter lemma with a stronger unproved premise is not a replacement for the
closed endpoint. Explicit typed clients and concrete assumption reports are
therefore part of the deliverable, not optional decoration.

\section{Further research questions}

The next questions are concrete enough to have measurable success criteria.

\begin{enumerate}
\item \textbf{Can the named and numerical bundle proofs share one normalization
      argument?} Both normalize requests, closure, and a private marker. An
      abstract matrix-satisfaction interface could expose \emph{semantic}
      freshness and update laws once, with named and decoded instances.
      Semantic freshness means that updating the private name does not change
      matrix truth. Requiring only syntactic absence of a free name would be
      too restrictive for the numerical code-zero case: its decoded tautology
      can mention name zero while its truth remains independent of that name.
      The existing \code{sat_dec_update} and the new \code{sat_update_fresh}
      already supply the two semantic interfaces. The target is a net
      deletion after including the interface, with the numerical diagonal case
      and irreflexivity assumptions still explicit.
\item \textbf{How far can the Rocq basis be reused beyond these 71 lemmas?}
      The unchanged meta-theory consumers already compile through the facades.
      A next step is to factor only the common contextual/deduction operations,
      while leaving the differing quantifier/vacuity hypotheses as parameters.
      The success criterion is paired source reduction plus preserved original
      signatures, not a claim that the two full calculi coincide.
\item \textbf{Can the direct lower bridge and the existing QHyp quotient adapter
      be adopted together?} They address different stages and should compose.
      The combined closure must be built afresh or through exact bound providers;
      source savings from two independent proposals cannot simply be added when
      they share imports or wrappers.
\item \textbf{Which proof-interface cuts improve build cost as well as source
      readability?} Measure clean and incremental builds with identical
      toolchains and resource settings. Keep generated certificate data,
      mathematical proof source, and compiled proof-object sizes separate.
\item \textbf{Can more cross-assistant architecture be aligned?} The Rocq direct
      seed construction suggested the Lean simplification; the Lean practice
      of shared named-calculus interfaces suggested the Rocq factoring. A
      theorem-interface correspondence could prevent future independent copies
      of essentially the same argument, without forcing both assistants to use
      identical encodings or proof tactics.
\item \textbf{Can the internal consistency statement be strengthened without
      duplicating the machine proof?} The internal PA endpoint is a distinct
      research problem. Isolate the additional formalized proof-translation and
      arithmetization obligations at the existing semantic/calculus boundary;
      do not rebuild a second physical runtime proof merely to change the
      metatheoretic presentation.
\end{enumerate}

\section{Sources and attribution}\label{sec:references}

All ProveIt URLs below refer to the immutable reviewed commit, not the moving
\code{main} branch. Long file paths are given so the report remains usable after
incorporation into the repository.

\begin{enumerate}[label={[\arabic*]}]
\item \href{https://github.com/VladimirReshetnikov/ProveIt/blob/6d9800d1ca67c52db0eb6e1c65df7b5cab9cdb95/Computability/BusyBeaver/Research/ZFC/candidate-k215/README.md}{K215 README}.
      Exact machine, theorem statements, evidence routes, and formal boundary.
\item \href{https://github.com/VladimirReshetnikov/ProveIt/blob/6d9800d1ca67c52db0eb6e1c65df7b5cab9cdb95/Computability/BusyBeaver/Research/ZFC/k225-formal/verification/k215/simplification/phase14/README.md}{Fourteenth K215 proof-simplification checkpoint}.
      Maintained source counts, checked scopes, and assumption profiles.
\item \href{https://github.com/VladimirReshetnikov/ProveIt/blob/6d9800d1ca67c52db0eb6e1c65df7b5cab9cdb95/Computability/BusyBeaver/Research/ZFC/k225-formal/verification/k215/simplification/proposals/README.md}{Proof-simplification proposals index}.
      Earlier source, runtime, scheduling, quotient, and interface proposals.
\item \href{https://github.com/VladimirReshetnikov/ProveIt/blob/6d9800d1ca67c52db0eb6e1c65df7b5cab9cdb95/Computability/BusyBeaver/Research/ZFC/k225-formal/verification/k215/simplification/proposals/sharing-interfaces/calculus/DIRECT_FULL_SEED.md}{Direct full seed proposal}.
      Prior uncompiled construction implemented by the Lean component here.
\item \href{https://github.com/VladimirReshetnikov/ProveIt/blob/6d9800d1ca67c52db0eb6e1c65df7b5cab9cdb95/Computability/BusyBeaver/Lean/BusyBeaver/ZFC/KEind222.lean}{\code{KEind222.lean}};
      \href{https://github.com/VladimirReshetnikov/ProveIt/blob/6d9800d1ca67c52db0eb6e1c65df7b5cab9cdb95/Computability/BusyBeaver/Lean/BusyBeaver/ZFC/KSeedModel222.lean}{\code{KSeedModel222.lean}};
      \href{https://github.com/VladimirReshetnikov/ProveIt/blob/6d9800d1ca67c52db0eb6e1c65df7b5cab9cdb95/Computability/BusyBeaver/Lean/BusyBeaver/ZFC/KIndisc222.lean}{\code{KIndisc222.lean}};
      \href{https://github.com/VladimirReshetnikov/ProveIt/blob/6d9800d1ca67c52db0eb6e1c65df7b5cab9cdb95/Computability/BusyBeaver/Lean/BusyBeaver/ZFC/KEndpoint.lean}{\code{KEndpoint.lean}}.
      Exact old lower-chain interfaces and proofs.
\item \href{https://github.com/VladimirReshetnikov/ProveIt/blob/6d9800d1ca67c52db0eb6e1c65df7b5cab9cdb95/Computability/BusyBeaver/Research/ZFC/k225-formal/rocq224/src/Complete222.v}{\code{rocq224/src/Complete222.v}}.
      Existing Rocq fresh-code construction and direct seed route.
\item \href{https://github.com/jeff87654/busy-beaver/tree/c132d7eae8158cd3a9850ea3eb7251158c3dd821}{Jeff Ketchersid's pinned busy-beaver source} and
      \href{https://github.com/elliotglazer/k230-rocq/tree/3142e4dd980f0be12f4abb636c4ae2528afc17eb}{Elliot Glazer's pinned K230 Rocq source}.
      The inherited named propositional derivations, source calculi, and
      set-theoretic construction retain their original authorship and terms.
\end{enumerate}

The abstract propositional derivation bodies are adaptations of the existing
Jeff Ketchersid proofs. The seed principle and lower construction are inherited
from the Glazer/K-family development. This packet's contribution is the
checked refactoring, concrete adapters, direct Lean implementation, integration
overlay, exact accounting, and review; it does not claim discovery of these
underlying logical principles. Notices and source identities accompany the code.

\appendix
\section{Archive map and reproducibility}
\begin{longtable}{@{}p{0.30\textwidth}p{0.65\textwidth}@{}}
\toprule
Path & Contents and role \\
\midrule
\endhead
\code{report/} & This article, its LaTeX inputs, and the compiled PDF. \\
\code{lean/} & Direct-seed source overlay, compatibility clients, reproduction
program, patch, source identities, and compilation evidence. \\
\code{rocq/} & Shared propositional core, two exact facades, 142-signature audit,
pinned input snapshots, deterministic producer, reproduction program, source
patches, notices, and compiler/kernel evidence. \\
\code{tools/closure_census.py} & Read-only Lean import and source-size census;
accepts an explicit checkout and repository-relative or module-root overlays. \\
\code{evidence/} & Repository pin, baseline and proposed closure measurements,
independent review, and delivery manifest. \\
\code{README.md} & Entry point, exact validation scope, integration instructions,
and instructions for rebuilding this article. \\
\bottomrule
\end{longtable}

To rebuild the PDF, run the following from \code{report/} in an environment
with pdfLaTeX and the listed standard packages:
\begin{lstlisting}
pdflatex -interaction=nonstopmode -halt-on-error main.tex
pdflatex -interaction=nonstopmode -halt-on-error main.tex
\end{lstlisting}
The reproduction programs do not require compiled proof objects from this
archive. Proof-object files and compiler installations are deliberately omitted;
the formal source, inputs or exact retrieval pins, and logs are provided.


\end{document}
